\documentclass[10pt,a4paper]{amsart}
\usepackage[margin=19mm]{geometry}
\usepackage{amsmath,amssymb,mathtools}
\usepackage[scr=rsfs,cal=euler]{mathalfa}
\usepackage{xfrac}
\usepackage[unicode,hidelinks,bookmarks=false]{hyperref}
\newcommand{\ie}{i.e.\ }
\newcommand{\cf}{cf.\ }
\newcommand{\resp}{respectively\ }
\newcommand{\Subsection}{\subsection}
\providecommand{\nobreakdash}{\nobreak\hbox{-}}
\setlength{\emergencystretch}{2em}
\multlinegap0pt
\usepackage{amssymb}
\usepackage{mathtools}
\usepackage{dsfont}
\usepackage{xcolor}
\usepackage{braket}
\usepackage{cancel}
\usepackage{bm}
\newtheorem{thm}{Theorem}
\newtheorem{lem}{Lemma}
\newcommand{\A}{\mathcal{A}}
\newcommand{\G}{\Gamma}
\newcommand{\LL}{\widetilde{\Li}}
\newcommand{\La}{\Lambda}
\newcommand{\Li}{\mathcal{L}}
\newcommand{\T}{T^\Lambda}
\newcommand{\TT}{\widetilde{T}^\Lambda}
\newcommand{\abs}[1]{\left \vert#1\right\vert}
\newcommand{\norm}[1]{\left\lVert#1\right\rVert}
\newcommand{\normm}[2]{\left\lVert#1\right\rVert_{#2}}
\title{Static features from mixing in short- and long-range Lindbladians: Markov property and correlations}
\author{Paul Rosa-Ruiz, Matteo Scandi, \'Angela Capel, \'Alvaro M. Alhambra}
\date{Source: arXiv:2606.28054v2 (CC BY 4.0)}
\begin{document}
\maketitle

\noindent\emph{Source and licence.} Static features from mixing in short- and long-range Lindbladians: Markov property and correlations. Version arXiv:2606.28054v2 (CC BY 4.0), distributed by arXiv under the Creative Commons Attribution 4.0 International licence, http://creativecommons.org/licenses/by/4.0/, as recorded in the arXiv licence metadata for this exact version and shown on the abstract page https://arxiv.org/abs/2606.28054v2. Full licence notice: \texttt{licenses/arxiv-2606.28054-CC-BY-4.0.txt}.\par\smallskip

\noindent\textit{Editorial scope.} Counted unit: the article's thm thm:clusteringBoundary (source0.tex lines 1413--1421). The article's complete original proof of it is delivered with it (source0.tex lines 1422--1644, one \texttt{proof} environment of 223 lines). No mathematics is written, repaired or re-derived: the statement and the proof are the article's own lines, in the article's own notation, and no retained line is substituted or re-flowed.  Retained uncounted as the source-local context the proof needs: lines 624--627; lines 1315--1323; lines 1407--1412; lines 2481--2486.  One declared adaptation: exactly 1 ancillary strings inside the retained spans are omitted (image-inclusion commands and, where the source repeats a label, the repeated label definitions) and no other line differs from the source; the figure environments, with their placement, captions and labels, are kept, and the package records the omitted strings in fidelity receipts bound to the affected bodies.  No retained line contains a citation, so the wrapper carries no bibliography and no bibliography entry.  The retained lines contain the article's own diagram code, which the wrapper typesets with the standard tikz packages; no image file is delivered and the package declares no asset dependency.  Editorial formatting only, in addition: an AMS \texttt{amsart} preamble that restates the article's own declarations and macros used by the retained lines, namely the environments \texttt{thm}, \texttt{lem} on separate counters, together with the standard packages those lines need.  The retained environments print in this document as Theorem 1, Lemma 1 in order of appearance, whereas the article prints the same items with the numbering of its own full document.  The complete native source of the article as distributed in the arXiv e-print for this exact version is bundled unchanged as \texttt{sources/203820/source0.tex}.
    \begin{equation}
        C_F:=\sup_{x,y\in \G}\sum_{z\in\G}\dfrac{F(d(x,z))F(d(z,y))}{F(d(x,y))}<\infty
        \label{eq:convolutionCF}
    \end{equation}

\begin{thm}\label{thm:localPerturbation}
    Let $(\Gamma,d)$ be a countable metric space equipped with an F-function $F$. Let $\Li=\{L_Z\}_{Z\in \mathcal{P}_0(\Gamma)}$ be a dissipative interaction with $\normm{\Li}{F}<\infty$ and primitive, satisfying local rapid mixing. Fix $A,C\in \mathcal{P}_0(\Gamma)$ with $R:=d(A,C)>0$. For any $\Lambda\in  \mathcal{P}_0(\Gamma)$ with $A\cup C \subset \Lambda$, $O_A\in \mathcal{A}_A$, a perturbation $Q_C\in \mathcal{CB}_0(C)$ and $\LL:=\Li+Q_C$ the perturbed \textit{Lindbladian}.
Let $\sigma$ be the unique fixed point of $\Li$ with convergence governed by $g(t)=e^{-\lambda t}$ for $\lambda >0$ and $\pi$ a fixed point of $\LL$.
Then,
    \begin{equation}\label{eq:diff_fixed_points_Lindbladians}
        \abs{(\sigma-\pi)(O_A)}\leq c\cdot \norm{O_A} \cdot k(\abs{A})\cdot F(R)^{\lambda/(\lambda +v)} \, ,
    \end{equation}
    where $v=\normm{\mathcal{L}}{F}\cdot C_F$, $k$ a polynomial function and a constant $c=1+1/(v\cdot C_F)$.
\end{thm}

\begin{figure}[!h]
    \centering
    
    \caption{Geometry considered for the MI: $A$ and $C$ and their frontier $\partial_{AC}$}
    \label{fig:MIgeometry}
\end{figure}

\begin{lem}\label{lem:technicalLemmaFfunction2}
    Let $\Gamma$ be a $D$-regular graph and let $\Li=\sum_{Z\in\mathcal{P}_0(\Gamma)}L_Z$ be a local interaction governed by the $F$-function $F_\alpha(r) :=1/(1+r)^\alpha;\,r\geq0$. Fix an $l>0$, then
    \begin{equation}
    \sup_{x\in\G}\sum\limits_{\substack{Z\in\mathcal{P}_0(\Gamma)\\ Z\ni x  \\ d_x(Z)= l}}\norm{L_Z}\leq \kappa \normm{\Li}{F} F_{\alpha -D+1}(l)
    \end{equation}
\end{lem}

\begin{thm}\label{thm:clusteringBoundary}
    Let $(\Gamma,d)$ be a countable metric space equipped with the F-function\[F_\alpha(r):=\dfrac{1}{(1+r)^\alpha},\]for $r\geq0$ and $\alpha>3D-1$. Let $\Li=\{L_Z\}_{Z\in \mathcal{P}_0(\Gamma)}$ be a local and primitive dissipative interaction with $\normm{\Li}{F}<\infty$, satisfying local rapid mixing. Fix $A,C\in \mathcal{P}_0(\Gamma)$ with $A\cap C =\emptyset$ (see Figure \ref{fig:MIgeometry}). For any $\Lambda\in  \mathcal{P}_0(\Gamma)$ with $A\cup C \subset \Lambda$, $O\in \mathcal{A}_{A\cup C}$, define the perturbation\footnote{We denote by $\partial_{AC}$ to the boundary separating $A$ and $C$ which is chosen equidistant to them (see Figure \ref{fig:MIgeometry}). With some abuse of notation, we denote $Z\cap \partial_{AC}\neq \emptyset$ to $Z$ that cross or touch the boundary $\partial_{AC}$.} $Q:=-\sum\limits_{\substack{Z\subset \Lambda \\ Z\cap \partial_{AC}\neq \emptyset}}L_Z\in CB_0(\partial_{AC})$ and $\LL:=\Li+ Q $ the perturbed \textit{Lindbladian}.
Let $\sigma\in\A_\La^*$ be the unique fixed point of $\Li$, which satisfies local rapid mixing ($\Li$ with a convergence governed by $g(t)=e^{-\lambda t}$ for $\lambda >0$), and let $\pi\in\A_\La^*$ a fixed point of $\LL$. Then,
    \begin{equation}
        \abs{(\sigma-\pi)(O)}\leq \norm{O}\cdot  k(\abs{A},\abs{C})\cdot F_{\alpha '}(R-1),
    \end{equation}
    where $\alpha'=\dfrac{\lambda(\alpha-2D)}{\lambda +v}$, $v=\normm{\Li}{F}\cdot  C_F$ the Lieb-Robinson speed, $k$ is a polynomial function, and $R=d(A,C)$. 

\end{thm}

\begin{proof}
    Notice that we can split the term $\abs{(\sigma-\pi)(O)}$ into two terms by using triangular inequality: the first term decays with time due to rapid mixing, the second term decays with $R$:
  \begin{align}
\abs{(\sigma-\pi)(O)} &= \abs{\sigma(O)-\pi(O)} \label{eq:diferenceStates} \\
&= \abs{\sigma(\T_t(O))-\pi(\TT_t(O))} \nonumber\\
&= \abs{(\sigma-\pi)(\T_t(O))-\pi((\TT_t-\T_t)(O))} \nonumber\\
&\le \abs{(\sigma-\pi)(\T_t(O))}+\abs{\pi((\TT_t-\T_t)(O))} \nonumber\\
&\le \underbrace{\abs{(\sigma-\pi)(\T_t(O))}}_{(1)}+\underbrace{\normm{(\TT_t-\T_t)(O)}{1}}_{(2)}. \nonumber
\end{align}


    For the first term, we use that $\Li$ satisfies local rapid mixing:
     \begin{equation*}
        (1):\quad \abs{(\sigma-\pi)(\T_t(O))} \leq p(\abs{A})\cdot   \norm{O} \cdot e^{-\lambda t}
     \end{equation*}
     For the second term, we will distinguish between two types of perturbation: those that intersect $A$ or $C$, and those that do not. This distinction motivates the following definitions:
     \begin{equation}
         \Omega:=\{ Z\subset \Lambda : Z\cap \partial_{AC}\neq \emptyset,Z\cap(A\cup C)= \emptyset\},
         \label{eq:Omega}
     \end{equation}
     \begin{equation}
         \Theta := \{ Z\subset \Lambda : Z\cap \partial_{AC}\neq \emptyset,Z\cap(A\cup C)\neq \emptyset\}.
         \label{eq:Theta}
     \end{equation}
     This decomposition separates perturbation terms according to whether
    they are spatially disjoint from the observable support $A\cup C$.
    The terms in $\Omega$ are at positive distance from $A\cup C$ and can
    therefore be controlled directly using the long-range Lieb–Robinson
    estimate. 
    
    In contrast, the terms in $\Theta$ intersect $A\cup C$ and cannot be
    treated purely by quasi-locality. This distinction is essential in the
    long-range setting, since interaction terms crossing $\partial_{AC}$
    may have arbitrarily large diameter.
    
    Using this decomposition, the perturbed Lindbladian can be written as
     \begin{equation}
         \LL=\Li -\sum_{Z'\in \Omega}L_{Z'}-\sum_{Z''\in \Theta}L_{Z''}
     \end{equation}
    we can rewrite the second term in the inequality above as:
     \begin{multline}
        (2):\normm{(\TT_t-\T_t)(O)}{1}\leq \int_0^t\underbrace{\sum_{Z'\in \Omega} \norm{L_{Z'}(\T_s(O))}}_{(4)}ds+\int_0^t\underbrace{\sum_{Z''\in \Theta} \norm{L_{Z''}(\T_s(O))}}_{(3)}ds
        \label{eq:(2)}
     \end{multline}
 Subsequently, we bound the term $(3)$ by applying Lemma \ref{lem:technicalLemmaFfunction2} and using the contractivity of $\T_t$.  
\begin{align}
(3):\;\sum_{Z''\in \Theta} \norm{L_{Z''}(\T_s(O))} &\le \sum_{Z''\in \Theta} \normm{L_{Z''}}{cb}\norm{\T_s(O)} \\
&\le \norm{O}\sum_{Z''\in \Theta} \normm{L_{Z''}}{cb} \nonumber\\
&\le \norm{O}\sum_{x\in A\cup C}\sum_{l\ge R/2}\underbrace{\sum_{\substack{Z''\subset \La \\ Z''\ni x \\ d_x(Z'')=l}}\norm{L_{Z''}}}_{\le \kappa \normm{\Li}{F} F_{\alpha-D+1}(l)} \nonumber\\
&\le \kappa \normm{\Li}{F}\norm{O}\abs{A\cup C}\underbrace{\sum_{l\ge R/2} F_{\alpha-D+1}(l)}_{\le \frac{1}{\alpha-D}F_{\alpha-D}(R/2-1)} \nonumber\\
&\le \norm{O}\underbrace{\kappa \normm{\Li}{F}\dfrac{1}{\alpha-D}\abs{A\cup C}}_{=:\tilde{c}'} F_{\alpha-D}(R/2-1). \nonumber
\end{align}
     
     For the terms of the Lindbladian that cross the boundary $\partial_{AC}$ but do not intersect $A\cup C$, we can use Theorem \ref{thm:localPerturbation}:
    \begin{align}
    (4):\;\sum_{Z'\in \Omega} & \norm{L_{Z'}(\T_s(O))}\nonumber \\
    &\nonumber \le \sum_{Z'\in \Omega}\left( \dfrac{\normm{L_{Z'}}{cb}\,\norm{O}}{C_F}\left( e^{\normm{\Li}{F} C_F s}-1\right)\sum_{x\in A\cup C}\sum_{y\in Z'}F_\alpha(d(x,y))\right) \\
    &= \left( \dfrac{\norm{O}}{C_F}\left( e^{\normm{\Li}{F} C_F s}-1\right)\underbrace{\sum_{Z'\in \Omega}\normm{L_{Z'}}{cb}\sum_{x\in A\cup C}\sum_{y\in Z'}F_\alpha(d(x,y))}_{(*)}\right).
    \label{eq:termsCrossingBoundary}
    \end{align}
    
    We will rewrite the term $(*)$ and split it into two terms (a) and (b) that are bounded independently:
    \begin{align}
(*):\;\sum_{Z'\in \Omega} \normm{L_{Z'}}{cb}&\sum_{x\in A\cup C}\sum_{y\in Z'}F_\alpha(d(x,y)) \nonumber\\ 
&\le \sum_{r\ge 1}\sum_{\substack{Z'\subset \La \\ d(A\cup C,Z')= r}}\normm{L_{Z'}}{cb}\sum_{x\in A\cup C}\sum_{y\in Z'}F_\alpha(d(x,y)) \label{eq:(*)split} \\
&= \underbrace{\sum_{r=1}^{R/2}\sum_{\substack{Z'\subset \La \\ d(A\cup C,Z')= r}}\normm{L_{Z'}}{cb}\sum_{x\in A\cup C}\sum_{y\in Z'}F_\alpha(d(x,y))}_{(a)} \nonumber\\
&\quad + \underbrace{\sum_{r\ge R/2+1}\sum_{\substack{Z'\subset \La \\ d(A\cup C,Z')= r}}\normm{L_{Z'}}{cb}\sum_{x\in A\cup C}\sum_{y\in Z'}F_\alpha(d(x,y))}_{(b)}. \nonumber
\end{align}
     

     
     Before continuing, note that
     \begin{equation}
         \sum_{x\in A\cup C}\sum_{y\in Z'}F_\alpha(d(x,y))\leq \text{min}\{\abs{A\cup C},\abs{Z'}\}\,\sup_{x\in \Lambda}\sum\limits_{\substack{y\in \Lambda \\ 
         d(x,y)\geq d(A\cup C,Z')
         }}   
         F_\alpha\left (d(x,y)\right).
     \end{equation}
    For the decay function $F_\alpha (r):=(1+r)^{-\alpha}$, a simple integration shows that for all $R\in \mathbb{N},\alpha > D$, there exists a constant $c > 0$ such that
    \begin{equation}
        \sup_{x\in \Lambda}\sum\limits_{\substack{y\in \Lambda \\ 
         d(x,y)\geq R
         }}   
         F_\alpha\left (d(x,y)\right)\leq \dfrac{c}{4}F_{\alpha-D}(R).
    \end{equation}
    In particular,
    \begin{align}
        \sum_{x\in A\cup C}\sum_{y\in Z'}F_\alpha(d(x,y))& \leq \dfrac{c}{4}\text{min}\{\abs{A\cup C},\abs{Z'}\}F_{\alpha-D}(d(A\cup C,Z')) \\ & \nonumber \leq\dfrac{c}{4}(\abs{A}+\abs{C})F_{\alpha-D}(d(A\cup C,Z')).
    \end{align}
    Going back to the Equation \eqref{eq:(*)split}. The first term $(a)$ corresponds to a sum over sets $Z\subset \La$ with diameter greater than $R/2-r$.
\begin{align}
(a):\;\sum_{r=1}^{R/2}&
\sum_{\substack{Z'\subset \La \\ d(A\cup C,Z')= r}}
\normm{L_{Z'}}{cb}
\underbrace{\sum_{x\in A\cup C}\sum_{y\in Z'}F_\alpha(d(x,y))}_{\le \dfrac{c}{4}\abs{A\cup C}F_{\alpha-D}(r)}
\\
&\le
\dfrac{c}{4}\abs{A\cup C}
\sum_{r=1}^{R/2}
F_{\alpha-D}(r)
\sum_{l\ge R/2-r}
\sum_{\substack{Z'\subset \La \\ d(A\cup C,Z')= r \\ d_x(Z')=l}}
\normm{L_{Z'}}{cb}
\nonumber\\
&\le
\dfrac{c}{4}\abs{A\cup C}
\sum_{r=1}^{R/2}
F_{\alpha-D}(r)
\sum_{l\ge R/2-r}
\sum_{\substack{x\in \La \\ d(x,\partial_{AC})= r}}
\sum_{\substack{x\in Z'\subset \La \\ d(A\cup C,Z')= r \\ d_x(Z')=l}}
\normm{L_{Z'}}{cb}
\nonumber\\
&\le
\dfrac{c}{4}\kappa \abs{A\cup C}
\sum_{r=1}^{R/2}
(1+r)^{D-1}
F_{\alpha-D}(r)
\sum_{l\ge R/2-r}
(\abs{\partial_{A}}+\abs{\partial_C})\underbrace{\sum_{\substack{x\in Z'\subset \La \\ d(A\cup C,Z')= r \\ d_x(Z')=l}}
\normm{L_{Z'}}{cb}}_{\le \kappa \normm{\Li}{F}F_{\alpha-D+1}(l)}
\nonumber\\
&\le
\dfrac{c}{4}\kappa^2 \normm{\Li}{F}
(\abs{\partial_A}+\abs{\partial_C})
\abs{A\cup C}
\sum_{r=1}^{R/2}
F_{\alpha-2D+1}(r) \underbrace{\sum_{l\ge R/2-r}F_{\alpha-D+1}(l)}_{\le \frac{1}{\alpha-D}F_{\alpha-D}(R/2-r-1)}
\nonumber\\
&\le
\dfrac{c}{4}\kappa^2 \normm{\Li}{F}
\frac{1}{\alpha-D}
(\abs{\partial_A}+\abs{\partial_C})
\abs{A\cup C}
\underbrace{\sum_{r=1}^{R/2}
F_{\alpha-2D+1}(r-1)
F_{\alpha-2D+1}(R/2-r-1)}_{\alpha-2D+1 >D \text{ needed for convolution}}
\nonumber\\
&\le
\dfrac{c}{4}\kappa^2 \normm{\Li}{F}
\frac{C_F}{\alpha-D}
(\abs{\partial_A}+\abs{\partial_C})
\abs{A\cup C}
F_{\alpha-2D+1}(R/2-1),
\nonumber
\end{align}

where we used the convolution property of the $F$-functions, see Equation \ref{eq:convolutionCF}. Remember that the convolution property is a property of $F$-functions, and $F_{alpha}$ is a $F$-function iff $\alpha >D$. That is the reason why the condition $\alpha -2D+1>D$ arises, because we need $F_{\alpha-2D+1}$ to be a $F$-function.

For the second term $(b)$, we use that $\norm{L_Z}\leq \normm{\Li}{F};\forall Z\subset \La $:
    \begin{align}
(b):\;\sum_{r\ge R/2+1}&\sum_{\substack{Z\subset \La \\ d(A\cup C,Z)= r}}\normm{L_{Z}}{cb}\sum_{x\in A\cup C}\sum_{y\in Z}F_\alpha(d(x,y))\\
&\le \dfrac{c}{4}\abs{A\cup C}\sum_{r\ge R/2+1}\sum_{\substack{Z\subset \La \\ d(A\cup C,Z)= r}}\normm{L_{Z}}{cb}F_{\alpha-D}(d(A\cup C,Z))
\nonumber\\
&\le \dfrac{c}{4}\abs{A\cup C}\sum_{r\ge R/2+1}F_{\alpha-D}(r)\sum_{l\ge 1}\sum_{\substack{x\in \La \\ d(x,\partial_{AC})= r}}\sum_{\substack{x\in Z\subset \La \\ d(A\cup C,Z)= r \\ d_x(Z)=l}}\normm{L_{Z}}{cb}
\nonumber\\
&\le \dfrac{c}{4}\abs{A\cup C}\sum_{r\ge R/2+1}F_{\alpha-D}(r)\,(\abs{\partial_{A}}+\abs{\partial_{C}})\,\kappa(1+r)^{D-1}\sum_{l\ge 1}\sum_{\substack{x\in Z\subset \La \\ d(A\cup C,Z)= r \\ d_x(Z)=l}}\normm{L_{Z}}{cb}
\nonumber\\
&\le \dfrac{c}{4}\kappa(\abs{\partial_{A}}+\abs{\partial_{C}})\abs{A\cup C}\sum_{r\ge R/2+1}F_{\alpha-2D+1}(r)\sum_{l\ge 1}\kappa \normm{\Li}{F}F_{\alpha-D+1}(l)
\nonumber\\
&\le \dfrac{c}{4}\kappa^2 \normm{\Li}{F}(\abs{\partial_{A}}+\abs{\partial_{C}})\abs{A\cup C}\sum_{r\ge R/2+1}F_{\alpha-2D+1}(r)\,\dfrac{1}{\alpha-D}\underbrace{F_{\alpha-D}(0)}_{=1}
\nonumber\\
&\le \dfrac{c}{4}\dfrac{1}{\alpha-D}\kappa^2 \normm{\Li}{F}(\abs{\partial_{A}}+\abs{\partial_{C}})\abs{A\cup C}\sum_{r\ge R/2+1}F_{\alpha-2D+1}(r)
\nonumber\\
&\le \dfrac{c}{4}\dfrac{1}{\alpha-D}\kappa^2 \normm{\Li}{F}(\abs{\partial_{A}}+\abs{\partial_{C}})\abs{A\cup C}\,\dfrac{1}{\alpha-2D}F_{\alpha-2D}(R/2)
\nonumber\\
&\le \dfrac{c}{4}\dfrac{1}{(\alpha-D)(\alpha-2D)}\kappa^2 \normm{\Li}{F}(\abs{\partial_{A}}+\abs{\partial_{C}})\abs{A\cup C}F_{\alpha-2D}(R/2).
\nonumber
\end{align}
    
    
    Going back to $(*)$ from Equation \eqref{eq:termsCrossingBoundary}:
    \begin{equation}
        (*):\sum_{Z'\in \Omega} \normm{L_{Z'}}{cb}\sum_{x\in A\cup C}\sum_{y\in Z'}F_\alpha(d(x,y))\leq \hat{c} F_{\alpha-2D}(R/2-1),
    \end{equation}
    where
    \begin{equation}
        \tilde{c}=\dfrac{c\cdot C_F}{4}\dfrac{(\alpha-2D +1/C_F)}{(\alpha-D )(\alpha-2D)}\kappa^2 \normm{\Li}{F} (\abs{\partial_{A}}+\abs{\partial_{C}})\abs{A\cup C}.
    \end{equation}

    We are now ready to bound the term $(4)$ of the Equation \eqref{eq:(2)}:
     
     \begin{equation}
         (4):\sum_{Z'\in \Omega} \norm{L_{Z'}(\T_s(O))}\leq\left ( \dfrac{\norm{O}}{C_F}\left ( e^{\normm{\Li}{F} C_F\cdot s}-1\right )\tilde{\tilde{c}}F_{\alpha-2D}(R/2-1)\right )
     \end{equation}
     with $
         \tilde{\tilde{c}}=\tilde{c}'+\tilde{c}
    $. Going back to the Equation \eqref{eq:(2)},
    \begin{multline}
        (2):\normm{(\TT_t-\T_t)(O)}{1}\\\leq  \dfrac{\tilde{\tilde{c}}\norm{O}}{C_F}\left ( \dfrac{e^{v\cdot t}-vt-1}{v} \right )F_{\alpha-2D}(R/2-1)ds        +\norm{O}\tilde{c}' F_{\alpha-D-1}(R/2-1)\cdot t
        \label{eq:clusteringTime}
     \end{multline}

     We are ready to bound Equation \eqref{eq:diferenceStates}:
    \begin{align}
|(&\sigma-\pi)(O)|\\
&\nonumber \le \norm{O}\Bigg[ p(\abs{A})e^{-\lambda t} + \dfrac{\tilde{\tilde{c}}}{C_F}\left(\dfrac{e^{vt}-vt-1}{v}\right)F_{\alpha-2D}(R/2-1) + \tilde{c}'\,F_{\alpha-D-1}(R/2-1)\,t \Bigg]
\\
&= \norm{O}\Big[ p(\abs{A})e^{-\lambda t} + c_1\,F_{\alpha-2D}(R/2-1)\,e^{vt} + c_2\,F_{\alpha-D-1}(R/2-1)\,t \Big] \nonumber\\
&\le \norm{O}\Big[ p(\abs{A})e^{-\lambda t} + \underbrace{(c_1+c_2\,c_3)}_{=:c}\,F_{\alpha-2D}(R/2-1)\,e^{vt} \Big]. \nonumber
\end{align}
     where $c_3:=\dfrac{1}{v\cdot e}=\max_{t\geq 0} t e^{-vt}$ is a constant so that $t\leq c_3 e^{vt}$. The other constants are
     \begin{equation}
         c_1:= \dfrac{\tilde{\tilde{c}}}{C_F \cdot v}=\kappa \dfrac{\normm{\Li}{F}}{C_Fv(\alpha-D)}\abs{A\cup C}+\dfrac{c\normm{\Li}{F}}{4vC_F}\dfrac{(\alpha-2D +1)}{(\alpha-D )(\alpha-2D)}\kappa^2 (\abs{\partial_{A}}+\abs{\partial_{C}})\abs{A\cup C}
     \end{equation}
     \begin{equation}
         c_2:= \tilde{c}'-\dfrac{\tilde{\tilde{c}}}{C_F}
     \end{equation}

    Now, using our freedom to choose any $t>0$, we look for a $t=t(R)$ so that\[p(\abs{A}) e^{-\lambda t}=\mathcal{C}v F_{\alpha-2D}(R/2-1) \cdot e^{vt}.\] Explicitly, we get that
    \begin{equation}
        t(R)=\log \left(  (c^{-1}p(\abs{A})F_{\alpha-2D}(R/2-1)^{-1})^{1/(\lambda+v)}  \right),
    \end{equation}
    and
    \begin{equation}
        c\cdot F_{\alpha-2D}(R/2-1) \cdot e^{vt}=c^{\lambda/(\lambda+v)}\cdot p(\abs{A})^{v/(\lambda +v)} F_{\alpha '}(R/2-1)
    \end{equation}
    with a long-range decay given by $\alpha'=\dfrac{\lambda(\alpha-2D)}{\lambda +v}$.
    We  conclude the proof by observing that:
    \begin{equation}
         \abs{(\sigma-\pi)(O)}\leq \norm{O} \underbrace{ 2^{1+\alpha'}c^{\lambda/(\lambda+v)} p^{v/(\lambda +v)} }_{=: k}F_{\alpha '}(R-1)\,.
    \end{equation}
\end{proof}

\end{document}
